\begin{theorem}[The sector coefficients along $\Delta$ from the arity-two equations]
\label{thm:c9-theta2}
Let $\hbar F_1^{(0)}$ contain the finite counterterm of Theorem \ref{thm:c4-finite}, and let the arity-one equations along $\Delta$, $P_b$ and $\xi\in\Sigma$ be solved on two legs by \eqref{eq:c9-F1Delta}, by $\hbar F_1^{(1)}(P_b)$ of Theorem \ref{thm:c9-P1} (i), with the sector coefficients $\theta_P,\theta_{P,\lambda},\theta_{P,V}$, and by $\hbar F_1^{(1)}(\xi)$ of Proposition \ref{prop:c11-oddsol} (i), with its free coefficients. The obstructions of arity two on two legs are the sums of the coefficients of $\varepsilon^0$ of the cycles with exactly one special edge and either two vertices $S^{(1)}(x)$, $S^{(1)}(y)$ or one vertex $S^{(2)}(x,y)$ on the cycle, the other vertices from $I$, as in Proposition \ref{prop:c9-structure} (ii).
\begin{itemize}
\item (i) With $[\Delta,P_b]=-2P_b$ (Remark \ref{rem:c2-dilatation}) and $S^{(2)}(\Delta,P_b)=0$ (Definition \ref{def:c2-S2}), the two-leg part of \eqref{eq:c9-arity2} along $(\Delta,P_b)$ is
\begin{equation}
Q\,F_1^{(2)}(\Delta,P_b)+2\,\text{units}\cdot\Big(\theta\,S^{(1)}(P_b)_\perp+\theta_\lambda\,S^{(1)}(P_b)_\lambda+\theta_V\,S^{(1)}(P_b)_V\Big)=\mathrm{Ob}^{(2)}(\Delta,P_b)+Q(\text{local}),
\label{eq:c9-theta2}
\end{equation}
with the sector parts $S^{(1)}(P_b)_d$ of Definition \ref{def:c11-kernels}; the coefficients along $P_b$ drop out. The classes $[S^{(1)}(P_b)_\perp]$, $[S^{(1)}(P_b)_\lambda]$, $[S^{(1)}(P_b)_V]$ are linearly independent in the cohomology of $Q$ on the quadratic local functionals of degree $-1$ and Weyl weight $5$, and
\[
\big[\hbar\,\mathrm{Ob}^{(2)}(\Delta,P_b)\big]=-\tfrac23\,\text{units}\cdot\Big([S^{(1)}(P_b)_\perp]+[S^{(1)}(P_b)_\lambda]+[S^{(1)}(P_b)_V]\Big),
\]
with no position part and no pole. Hence the arity-two equation along $(\Delta,P_b)$ is solvable on two legs exactly when $\theta=\theta_\lambda=\theta_V=-\tfrac13$.
\item (ii) With $[\Delta,\xi]=-\xi$, for $\xi$ of either chirality, the two-leg part of \eqref{eq:c9-arity2} along $(\Delta,\xi)$ is solvable exactly when $\theta=\theta_\lambda=\theta_V=-\tfrac13$, for every value of the free coefficients along $\xi$. The obstruction is
\[
\hbar\,\mathrm{Ob}^{(2)}(\Delta,\xi)=\text{units}\cdot\Big[8(T_1+T_2)-\tfrac{26}3(T_3+T_4)+\log\tfrac43\,Q(\psi_1+\psi_2)\Big],\qquad \psi_1=\int x^\mu\big(\Gamma(\xi,\lambda)^\mu,c^*\big),\quad \psi_2=\int x^\mu\big(A^*_\mu,(\xi,\lambda^*)\big),
\]
with $T_1,\dots,T_4$ the terms of Table \ref{tab:c11-terms} and $\Gamma(\xi,\lambda)^\mu$ the component of $\Gamma(\xi,\lambda)|_V$ along $e_\mu$, and at these values $\hbar F_1^{(2)}(\Delta,\xi)=\log\tfrac43\,\text{units}\cdot(\psi_1+\psi_2)$ plus $Q$-closed functionals.
\item (iii) At $\theta=\theta_\lambda=\theta_V=-\tfrac13$ the quadratic part of $\hbar F_1^{(1)}(\Delta)$ in the sector $(A|_{W_{\mathbb R}^\perp},A^*|_{W_{\mathbb R}^\perp})$ is $\big[\tfrac{20}3\int(A|_{W_{\mathbb R}^\perp},A^*|_{W_{\mathbb R}^\perp})-\tfrac23\int(E(A|_{W_{\mathbb R}^\perp}),A^*|_{W_{\mathbb R}^\perp})\big]$ units, and the class of $\hbar F_1^{(1)}(\Delta)$ minus the translation-invariant solution is $-\tfrac13$ units times $[S^{(1)}(\Delta)]$: the uniform coefficient along $\Delta$ is $-\tfrac13$, independently of the coefficients along $P_b$, $\Sigma$ and $\Sigma^*$ and of the field redefinition $\phi^{(0)}=c\,\text{units}\cdot S^{(1)}(\Delta)$, which shifts the coefficient of $S^{(1)}(x)$ in $\hbar F_1^{(1)}(x)$ by $c\lambda_x$ units, $\lambda_\Delta=0$.
\end{itemize}
\end{theorem}

\begin{proof}
(i) In \eqref{eq:c9-arity2} at $(\Delta,P_b)$ the term $\{F_1^{(0)},S^{(2)}(\Delta,P_b)\}$ vanishes and the coefficients enter through $F_1^{(1)}(-2P_b)$, $\{F_1^{(1)}(\Delta),S^{(1)}(P_b)\}$ and $\{S^{(1)}(\Delta),F_1^{(1)}(P_b)\}$. The sector parts of $S^{(1)}(\Delta)$ and of $S^{(1)}(P_b)$ in different sectors have no common legs, and within one sector $\{S^{(1)}(\Delta)_d,S^{(1)}(P_b)_d\}=v.S^{(1)}(P_b)_d=2S^{(1)}(P_b)_d$, $S^{(1)}(P_b)_d$ having Weyl weight $5$ (Lemma \ref{lem:c9-SDelta} (i)). The terms are therefore $2\theta_d\,S^{(1)}(P_b)_d$ units from $\{F_1^{(1)}(\Delta),S^{(1)}(P_b)\}$, $2\theta_{P,d}\,S^{(1)}(P_b)_d$ units from $\{S^{(1)}(\Delta),F_1^{(1)}(P_b)\}$ and $-2\theta_{P,d}\,S^{(1)}(P_b)_d$ units from $F_1^{(1)}(-2P_b)$. The translation-invariant solution of \eqref{eq:c9-F1Delta} is a combination of field-counting functionals, whose brackets with $S^{(1)}(P_b)$ vanish, and the $\kappa$-directions are $Q$-exact, with $Q$-exact brackets. Independence: $Q$ preserves the three sectors $(A|_{W_{\mathbb R}^\perp},A^*|_{W_{\mathbb R}^\perp})$, $(\lambda,\lambda^*)$ and $(A|_V,A^*|_V)\oplus(c,c^*)$ of the quadratic functionals of degree $-1$. In the first there is no functional of degree $-2$, $\int f(x)(A^*|_{W_{\mathbb R}^\perp},A^*|_{W_{\mathbb R}^\perp})$ vanishing for the odd $A^*$ and the symmetric fibre pairing, so $[S^{(1)}(P_b)_\perp]\neq0$. In the second the translation-invariant image is spanned by $Q\int(\rho(b)\lambda^*,\lambda^*)=2\int(\partial_b\lambda,\lambda^*)+2\int(\rho(b\wedge\partial)\lambda,\lambda^*)$, not a multiple of $S^{(1)}(P_b)_\lambda$. In the third it is spanned by $Q\int(\iota_bA|_V,c^*)$, which differs from $S^{(1)}(P_b)_V$ by the nonzero $\int(\iota_bdA|_V,A^*|_V)$ up to sign. The class of the obstruction is the sum of the coefficients of $\varepsilon^0$ of the cycles with the vertices $S^{(1)}(\Delta)$, $S^{(1)}(P_b)$ and one special edge, computed exactly on complete bases of the two-leg functionals with these quantum numbers; the residual of Lemma \ref{lem:c9-sform} is then $(-\tfrac23-2\theta_d)$ units times $[S^{(1)}(P_b)_d]$ in the sector $d$. (ii) The same computation along $(\Delta,\xi)$, with $S^{(2)}(\Delta,\xi)=0$: the coefficients along $\xi$ enter through $F_1^{(1)}(-\xi)$ and $\{S^{(1)}(\Delta),F_1^{(1)}(\xi)\}$, which cancel, $S^{(1)}(\xi)$ having Weyl weight $\tfrac92$, and the sector coefficients along $\Delta$ through $\{F_1^{(1)}(\Delta),S^{(1)}(\xi)\}$, whose sector parts act on the legs of $T_1,\dots,T_4$; the cohomology on two legs at Weyl weight $\tfrac92$ and position degree at most one is computed on complete bases, and the residual class vanishes exactly at the stated values. (iii) Insert $\theta=-\tfrac13$ into \eqref{eq:c9-Fperp}. The field redefinition generated by $\phi^{(0)}$ changes $F_1^{(1)}(x)$ by $\{\phi^{(0)},S^{(1)}(x)\}$, which is $c\lambda_x$ units times $S^{(1)}(x)$ modulo exact terms (Lemma \ref{lem:c2-components}, $S^{(2)}(\Delta,x)=0$).
\end{proof}
